\documentclass[a4paper, 12pt, UTF8]{article}

% preamble
% package
\usepackage[BoldFont]{xecjk}
\setCJKmainfont[BoldFont=SimHei]{SimSun}
\setCJKmonofont{SimSum}
\usepackage{geometry}
\usepackage{graphicx}
\usepackage{float}
\usepackage{caption}
\usepackage{enumerate}


% command
\graphicspath{{pic/}}
\geometry{
    textheight=230mm
}

\begin{document}

\title{\Huge E2Ch Week 1\_2}
\author{\Large 
        1931991 Ren Jiaping \\[12pt]
        Tongji University \\[12pt]
        College of Surveying and GeoInfoMatics}
\date{2019-09-11}
\maketitle

\thispagestyle{empty}

\newpage
\pagenumbering{Roman}
\tableofcontents
\addtocontents{toc}{\protect\vspace{10pt}}

\newpage
\pagenumbering{arabic}
\part{Homework}

\newpage
\section{Text(Part)}

\begin{bfseries}
    \Large
    Grand openings
    \paragraph*{}
    \large
    Changes that will bring scientific discovery more freely into the public domain are happening. About time too.
\end{bfseries}

\paragraph*{}
    IN 2001 a meeting on scientific publishing held in Budapest by what was then called the Open Society Institute (now the Open Society Foundation) coined the phrase “open access”. The gathering official statement asked the world to “share the learning of the rich with the poor and the poor with the rich, make this literature as useful as it can be, and lay the foundation for uniting humanity in a common intellectual conversation” --- in other words, to make scientific papers free to users.

\paragraph*{}
    A noble aspiration, but one which cynics might have thought had little chance of coming to fruition. The rich, they would observe, include academic publishers, who have enjoyed three centuries of dominion over the dissemination of scientific work and who often have profit margins approaching 40\%. They had every incentive to scupper change.

\paragraph*{}
    Cynicism, however, is not always correct. The open-access movement which the meeting helped spawn now looks unstoppable. All seven of Britain’s research councils, for example, now require that the results of the work they pay for are open-access in some way. So does the Well-come Trust, a British charity whose medical-research budget exceeds that of many scientifically successful countries. And by 2016 every penny of public money given to British universities by the government will carry the same requirement.

\paragraph*{}
    Elsewhere, the story is the same. In 2013, after years of wrangling in America’s Congress, the White House stepped in to require federal agencies that spend more than \$100m a year on research to publish the results where they can be read for free. Countless universities, societies and funding bodies in other countries have similar requirements.

\newpage
\section{First Edition}
\subsection{Translation}
\begin{figure}[H]
    \centering
    \includegraphics[width = \textwidth]{original_work_01.jpg}
    \caption{First\_Edition\_01}
    \label{Fig:1}
\end{figure}

\begin{figure}[htbp]
    \centering
    \includegraphics[width = \textwidth]{original_work_02.jpg}
    \caption{First\_Edition\_02}
    \label{Fig:2}
\end{figure}

\subsection{Review}
Here records the mistakes of my translation.
\begin{enumerate}[\hspace{0.5cm} 1.]
    \item {\bfseries public domain} 公共领域
          \begin{itemize}
              \item \emph{MyTranslation:} 公共视线
              \item I translated it into 公共领域 at first time, then I modified. So I was really unsure about how to do this translation for lacking reading exprience. 
          \end{itemize}
    \item {\bfseries About time too.} It's about time to make this change too.
          \begin{itemize}
              \item \emph{MyTranslation:} It was translated as faster compared with \lq more freely\rq\ in the context above. Or \lq time that ... is happening\rq\ according to the structure of \lq changes that ... are happening\rq.
              \item Refer to the context. MORE READING, MORE THINKING!!! 
          \end{itemize}
    \item {\bfseries the learning of ...} 知识
          \begin{itemize}
              \item \emph{MyTranslation:} 受教育
              \item Too much abstraction and also panic about the time limit.
          \end{itemize}
    \item {\bfseries literature} 论文, 知识
          \begin{itemize}
              \item \emph{MyTranslation:} 文学
              \item Refer to context, do not and never translate directly!
          \end{itemize}
    \item {\bfseries common intellectual conversation} 学术交流
          \begin{itemize}
              \item \emph{MyTranslation:} 共同智力会谈 
              \item So it's a stupid direct translate again. Refer to the context!!! and THINK!!! Let the words swirl around your rooted mind!!!
          \end{itemize}
    \item {\bfseries humanity} 人类
          \begin{itemize}
              \item \emph{MyTranslation:} 人性
              \item Panic. I don't even know how to explain it. 
          \end{itemize}
    \item {\bfseries margins} 顶, 边
          \begin{itemize}
              \item \emph{MyTranslation:} 保证金 
              \item I know margin has a meaning of \lq 顶, 边\rq\ in Chinese. 
          \end{itemize}
    \item {\bfseries spwan} 产生
          \begin{itemize}
              \item \emph{MyTranslations:} 传播
              \item Hard to explain it from my twisted brain.
          \end{itemize}
\end{enumerate}

\newpage
\section{Second Edition}
\subsection{Translation}
\subsection{Review}

\newpage
\part{Complementary Materials}

\newpage
\section{Text(entire)}
\begin{bfseries}
    \Large
    Grand openings
    \paragraph*{}
    \large
    Changes that will bring scientific discovery more freely into the public domain are happening. About time too.
\end{bfseries}

\paragraph*{}
    IN 2001 a meeting on scientific publishing held in Budapest by what was then called the Open Society Institute (now the Open Society Foundation) coined the phrase “open access”. The gathering official statement asked the world to “share the learning of the rich with the poor and the poor with the rich, make this literature as useful as it can be, and lay the foundation for uniting humanity in a common intellectual conversation” --- in other words, to make scientific papers free to users.

\paragraph*{}
    A noble aspiration, but one which cynics might have thought had little chance of coming to fruition. The rich, they would observe, include academic publishers, who have enjoyed three centuries of dominion over the dissemination of scientific work and who often have profit margins approaching 40\%. They had every incentive to scupper change.

\paragraph*{}
    Cynicism, however, is not always correct. The open-access movement which the meeting helped spawn now looks unstoppable. All seven of Britain’s research councils, for example, now require that the results of the work they pay for are open-access in some way. So does the Well-come Trust, a British charity whose medical-research budget exceeds that of many scientifically successful countries. And by 2016 every penny of public money given to British universities by the government will carry the same requirement.

\paragraph*{}
    Elsewhere, the story is the same. In 2013, after years of wrangling in America’s Congress, the White House stepped in to require federal agencies that spend more than \$100m a year on research to publish the results where they can be read for free. Countless universities, societies and funding bodies in other countries have similar requirements.

\paragraph*{}
    Publishers, though they have often dragged their feet, are adjusting. This week the oldest, the Royal Society, and arguably the most prestigious, Nature Publishing Group (NPG)—both based in London—joined in. Each will now publish a journal that readers do not have to pay to look at.

\paragraph*{}
    \begin{bfseries}
        \large
        Who will publish, and who perish?
    \end{bfseries}


    Open-access publishing actually began in 2000, a year before the Budapest meeting, with the launch, in Britain, of BioMed Central, and in America, of the Public Library of Science (PLOS). The open-access business model shifts the cost of producing the journal from subscribers, such as university libraries, to researchers themselves, who pay an article-processing charge (APC) to appear in print or its electronic equivalent. Either way, the taxpayer picks up the bill in the end. But open publication makes research more widely accessible, which is a public good in its own right.

\paragraph*{}
    One problem open access brings is a shift in incentives. More published articles means more revenue from processing charges. Rejection rates at high-end paid-for journals often exceed 90\%. A commercial open-access publisher (which PLOS, a charity, is not) might be tempted to publish anything that came his way, in order to pocket the APC—and many have, giving open access a fly-by-night feel to some academics.

\paragraph*{}
    For example, in a survey conducted by NPG, of 27,000 authors of papers in its journals, 44\% expressed some concern about the quality of open-access publishing and more than a third agreed with the notion that it was associated with less prestige. Yet those perceptions may be misguided. A study of Nature Communications (which was, until this week’s announcement, a hybrid between open access and traditional subscription, but is now pure open access) shows that its open-access papers enjoyed a slightly higher number of citations and significantly more downloads and online views than their non-open-access brethren.

\paragraph*{}
    To foster prestige, champions of open access have launched efforts such as eLife, an online journal with an array of academic bigwigs at its helm. They hope to create a top-tier publication by sheer weight of bigwiggery. But eLife and its kind cannot force the change alone, because publishing with the requisite due diligence is not cheap.

\paragraph*{}
    The Public Library of Science has found a way out of this by copying the more-the-merrier approach, but in a controlled way. Until 2006 it was a producer of high-impact but loss-making publications. Then it started PLOS ONE, a different kind of journal altogether. Instead of acting as an arbiter of the importance of scientific work, PLOS ONE claims only to ensure that articles are scientifically sound. With less effort going into peer review, PLOS ONE publishes many more papers (in 2013, it carried more than 31,000 articles, 36 times as many as the next-most-prolific PLOS journal) while simultaneously charging less for them and becoming a cash-cow that helps pay for the rest of the outfit. The Royal Society hopes to recapitulate this idea with its new offering, Royal Society Open Science.

\paragraph*{}
    \begin{bfseries}
        \large
        Free hits
    \end{bfseries}

    Whether the experience of Nature Communications will overcome researchers’ misgivings remains to be seen. Despite the Wellcome Trust’s requirement that its grantees publish in open-access journals, only 70\% do so. To ensure compliance, the trust has had to introduce punitive measures, such as withholding money.

\paragraph*{}
    Some researchers just don’t care, though. A survey by Taylor \& Francis, a publishing firm, asked American and British scientists if they had published under open-access policies; 44\% and 32\%, respectively, did not know. More than half responded they did not know if they would in future. But many do know, and resist.

\paragraph*{}
    The point about prestige is not mere snobbery. The league table of journals is as finely graded as that of football, and, at the moment, has far less scope for promotion and demotion. Grant-awarding bodies and appointment committees know this and, in a wonderful display of doublethink, promote open access but also promote those who eschew it by publishing in top-notch, non-open journals

\paragraph*{}
    Only when that changes will open access’s victory be complete. This could happen either by new open-access journals acquiring the necessary kudos, or by old ones, seeing the game is up, becoming open access themselves. Though Nature Communications is a successful and well-regarded publication, it is not NPG’s top product. And Royal Society Open Science is untested. At the moment, then, both the Royal Society and NPG seem to be hedging their bets. When the society’s Proceedings, and NPG’s eponymous flagship, Nature, are both free for anyone to read, then open access’s partisans really will be able to declare victory and go home.

\section{Translation}
\section{Summary}
\end{document}